\documentclass[
aspectratio=169,
%	english,
design=2023,
german,
accentcolor=TUDa-2c,
% headrule=true, %OPTIONAL: Hinzufügen für farbigen Streifen (in diesem Fall TUDa-2c)
author-on-title,
logofile=images/tuda_logo.pdf
]{tudabeamer} 

%%%%%%%%%%%%%%%%%%%%%%%%%%%%
% Download des TU-Logos
%%%%%%%%%%%%%%%%%%%%%%%%%%%%
% https://download.hrz.tu-darmstadt.de/protected/CE/TUDa_LaTeX/tuda_logo.pdf
% The path to the logo can be given a "logofile" above

\usepackage[ngerman]{babel}
\usepackage[babel]{csquotes}
\usepackage{media9}
\usepackage{animate}
\usepackage[overlay]{textpos}
\usepackage{tuda-pgfplots}

\usepackage[utf8]{inputenc}
\usepackage[ngerman]{babel}
\usepackage{hyperref}
\usepackage{graphicx}
\usepackage{subcaption}

\usepackage{amsmath, amsthm, amssymb}
\usepackage[squaren]{SIunits}

\title{Differentiation und Integration}
\subtitle{\small Kurs: Wissenschaftliches Arbeiten im CE}
\author[Karakoc | Rosenhagen | Tanase]{Aykut Karakoc, Christian Rosenhagen, Bogdan Tanase}
\department{TU Darmstadt | Computational Engineering}
\institute{WACE -- SoSe 2024}

\titlegraphic*{\includegraphics[]{images/tuda_home2.jpg}}
\partnerlogo{\textcolor{accentcolor}{\includegraphics[height=0.7cm]{images/ce_logo_black.pdf}}}
%\date{\today}
\date{10. Juli 2024}

\newcommand{\pdfmovie}[4]{\href{run:#1}{\framebox{\parbox[c][#3][c]{#2}{\center #4}}}}




\begin{document}
	\maketitle
	
	\setbeamertemplate{title in head/foot}{Normalerweise auch Titel + Authoren}
	\begin{frame}
		\frametitle{Inhalt des Vortrags}
		\tableofcontents%[currentsection,subsectionstyle=show/show/hide]
	\end{frame}
	
	\section{Section 1}
	\subsection{Subsection 1.1}
	\subsection{Subsection 1.2}
	\section{Section 2}
	
%	\begin{frame}{Gliederung}
%		\begin{itemize}[<+->]
%			\item Thema 1
%			\item Thema2
%		\end{itemize}
%	\end{frame}
	
	
	\setbeamertemplate{title in head/foot}{Bald wird nichts mehr im header sein}
	\begin{frame}
		\frametitle{Vorlage 1}
	\end{frame}

	\setbeamertemplate{title in head/foot}{}
	\begin{frame}
		\frametitle{Vorlage 2}
	\end{frame}

	\begin{frame} 
	\frametitle{Quellen (noch mit Beispiel)}
		\begin{thebibliography}{9}
			\bibitem{Pan2019} \textcolor{darkgray}{P. Pan und T. Bell, \emph{Theorie über farbenfrohe Pferde}, 2019, Wunschverlag.}
		\end{thebibliography}
	\end{frame}


\end{document}